\documentclass[10pt,a4paper]{amsart}
\usepackage[margin=19mm]{geometry}
\usepackage{amsmath,amssymb,mathtools}
\usepackage[scr=rsfs,cal=euler]{mathalfa}
\usepackage{xfrac}
\usepackage[unicode,hidelinks,bookmarks=false]{hyperref}
\newcommand{\ie}{i.e.\ }
\newcommand{\cf}{cf.\ }
\newcommand{\resp}{respectively\ }
\newcommand{\Subsection}{\subsection}
\providecommand{\nobreakdash}{\nobreak\hbox{-}}
\setlength{\emergencystretch}{2em}
\multlinegap0pt
\usepackage{amssymb}
\newtheorem{lemma}{Lemma}
\newcommand{\bm}[1]{\boldsymbol{#1}}
\newcommand{\mc}[1]{\mathcal{#1}}
\title{A Correction Function-based KFBI Method for Brinkman Interface Problems}
\author{Han Zhou, Wenjun Ying}
\date{Source: arXiv:2604.15509v1 (CC BY 4.0)}
\begin{document}
\maketitle

\noindent\emph{Source and licence.} A Correction Function-based KFBI Method for Brinkman Interface Problems. Version arXiv:2604.15509v1 (CC BY 4.0), distributed by arXiv under the Creative Commons Attribution 4.0 International licence, http://creativecommons.org/licenses/by/4.0/, as recorded in the arXiv licence metadata for this exact version and shown on the abstract page https://arxiv.org/abs/2604.15509v1. Full licence notice: \texttt{licenses/arxiv-2604.15509-CC-BY-4.0.txt}.\par\smallskip

\noindent\textit{Editorial scope.} Counted unit: the article's lemma lemma (source0.tex lines 693--695). The article's complete original proof of it is delivered with it (source0.tex lines 696--754, one \texttt{proof} environment of 59 lines). No mathematics is written, repaired or re-derived: the statement and the proof are the article's own lines, in the article's own notation, and no retained line is substituted or re-flowed.  No further source-local context is needed: every cross-reference of every retained line resolves inside the two delivered spans.  Nothing was omitted from any retained span, so the package carries no omission record.  No retained line contains a citation, so the wrapper carries no bibliography and no bibliography entry.  No retained line contains a figure environment, an image-inclusion command or drawing code, so no image is delivered and the package declares no asset dependency.  Editorial formatting only, in addition: an AMS \texttt{amsart} preamble that restates the article's own declarations and macros used by the retained lines, namely the environments \texttt{lemma} on separate counters, together with the standard packages those lines need.  The retained environments print in this document as Lemma 1 in order of appearance, whereas the article prints the same items with the numbering of its own full document.  The complete native source of the article as distributed in the arXiv e-print for this exact version is bundled unchanged as \texttt{sources/198610/source0.tex}.
\begin{lemma}
    Let $\Gamma\cap \mc O_i$ be a straight line segment. Then $\bm M$ is independent of $R$ and is invertible.
\end{lemma}

\begin{proof}
We can write $\Gamma\cap \mc O_i$ as
\begin{equation}
    \Gamma\cap \mc O_i = \{(x,y):x\in (-R, R), y = 0\}.
\end{equation}
Then, the collocation points are given by $(\eta_k^{(1)} R, 0)\in \Gamma\cap \mc O_i$ with $k = 1,2,3$ for the Dirichlet boundary condition, $(\eta_k^{(2)} R, 0)\in \Gamma\cap \mc O_i$ with $k = 1,2$ for the Neumann boundary condition. The collocation point for the Brinkman equation is chosen as $(0, 0)\in \mc O_i$ and the divergence free condition $(\eta_1^{(4)}, 0), (\eta_2^{(4)}, 0), (0, \eta_3^{(4)})$ with $\eta_3^{(4)}\neq 0$.


To show that $\bm M$ is invertible, we only need to demonstrate that the homogeneous system $\bm M \bm c = \bm 0$ has only a zero solution.
The Dirichlet boundary condition simplifies to
\begin{equation}
    \bm c_1 + \eta_k^{(1)} \bm c_2 +  (\eta_k^{(1)} )^2\bm c_4 = \bm 0, \quad k = 1,2,3.
\end{equation}
Since $\eta_k^{(1)}$ are distinct, this is equivalent to the polynomial interpolation in 1D and the coefficient matrix is the Vandermonde matrix, which is invertible. Then, we have $\bm c_1 = \bm c_2 = \bm c_4 = \bm 0$. 
The collocation equation for the Neumann boundary conditions and the equation for the divergence-free condition at $\bm q_1^{(4)}, \bm q_2^{(4)}$ lead to
\begin{equation}
    \begin{aligned}
        \bm c_3 + \eta_k^{(2)}\bm c_6 + c_3^{(2)} 
        \begin{bmatrix}
            0\\
            1
        \end{bmatrix}
        + c_6^{(2)}
        \begin{bmatrix}
            0\\
            \eta_k^{(2)}
        \end{bmatrix}
        - d_1
        \begin{bmatrix}
            0\\
            1
        \end{bmatrix}
        -d_2
        \begin{bmatrix}
            0\\
            \eta_k^{(2)}
        \end{bmatrix} &= \bm 0, \\
    c_3^{(2)} + \eta_k^{(2)} c_6^{(2)} &= 0.
    \end{aligned}
\end{equation}
Since $\eta_1^{(2)} \neq \eta_2^{(2)}$ and $\eta_1^{(4)} \neq \eta_2^{(4)}$, we obtain
\begin{equation}
    \bm c_3 = \bm c_4=\bm 0,\quad d_1 = d_2 = 0.
\end{equation}
Considering the equation at $\bm q^{(3)}$ and $\bm q_3^{(4)}$, we have
\begin{equation}
\begin{aligned}
    2\bm c_5 - d_3 
    \begin{bmatrix}
        0 \\
        1
    \end{bmatrix} = \bm 0,\quad
    2\eta_3^{(4)} c_5^{(2)} = 0.    
\end{aligned}
\end{equation}
Since $\eta_3^{(4)}\neq 0$. We conclude $\bm c_5 = \bm 0$ and $d_3 = 0$. 
Hence, we have zero solution $\bm c =\bm 0$.
The above calculation also shows that $\bm M$ is independent of $R$.
\end{proof}

\end{document}
